\subsection{The smallest odd-prime case reduces to two levels}

\begin{lemma}[Finite inertia in a height-two formal group at $3$]
\label{lem:B-RD31-finite}
Let $G/K$ have potentially good reduction to a one-dimensional
formal BT group of height two, with residue characteristic $3$.
After passage to the completed maximal unramified extension, it
has good reduction over a Galois extension of degree at most $12$.
Its faithful special inertia image is either cyclic of order
dividing $8$, or has a normal subgroup of order $3$ and order at
most $12$; in the latter case its tame quotient has order $1,2$,
or $4$. Two such groups have simultaneous good reduction over an
extension of ramification degree at most $36$.
\end{lemma}
\begin{proof}
We work with algebraically closed residue field. Lazard's
classification theorem \cite[Theorem~IV, pp.270--271]{Lazard1955Formal}
identifies the special formal group with the height-two Honda
group. Its rational Dieudonne module has a basis in which
\[
 F=\begin{pmatrix}0&3\\1&0\end{pmatrix}\sigma.
\]
This is also obtained directly from the Dieudonne module of the
supersingular elliptic curve $Y^2=X^3-X$; a change of basis over
$W(\bar k)[1/3]$ absorbs the sign in $F^2=-3$.
The endomorphisms commuting with $F$ are precisely
\[
 A(a,c)=\begin{pmatrix}a&3\sigma(c)\\c&\sigma(a)\end{pmatrix},
       \qquad a,c\in\mathbf Q_{3^2}.
                                                        \tag{B-RD31.1}
\]
Here $\mathbf Q_{3^2}=W(\mathbf F_9)[1/3]$ and $\sigma$ is its
nontrivial automorphism. The determinant is
$N(a)-3N(c)$, whose two nonzero terms have different parity of
$3$-adic valuation. Thus this algebra is a quaternion division
algebra. Its valuation ring consists of the matrices with
$a,c\in W(\mathbf F_9)$; its residue field is $\mathbf F_9$, and
the kernel of reduction on units is pro-$3$. These facts also
follow by writing the algebra as
$\mathbf Q_{3^2}\langle\Pi\rangle$, with
$\Pi^2=3$ and $\Pi a=\sigma(a)\Pi$.

Let $H$ be any finite subgroup of its units. Its reduction has
cyclic image in $\mathbf F_9^\times$. Its kernel $P$ is a finite
$3$-group. If $P\ne1$, choose a central element $z$ of order $3$.
Its centralizer is the quadratic field $\mathbf Q_3(z)$, so all
of $P$ lies in that field. Its finite $3$-power roots of unity
form $\mu_3$: an element of order $9$ would have minimal
polynomial of degree $6$, whereas every element of the
quaternion algebra satisfies a quadratic polynomial. Hence
$P=C_3$. The centralizer of $P$ in $H$ lies in the roots of unity
of $\mathbf Q_3(\zeta_3)$, namely $\mu_6$. Conjugation embeds
$H/C_H(P)$ in $\operatorname{Aut}(C_3)=C_2$. Consequently
$|H|\leq12$. Its tame quotient has order dividing $8$ and at
most $4$, as asserted. If $P=1$, reduction embeds $H$ in the
cyclic group of order $8$.

The faithful special inertia image gives a good-reduction
extension of that degree by exactly the descent and lattice
argument of Lemma~\ref{lem:B-RD17-degree}. In particular, the
good model remains a one-dimensional formal group, since its
special isocrystal remains of slope $1/2$.

For two groups, their joint wild image embeds in $P_1\times P_2$,
and the joint tame quotient is cyclic. If both wild images are
nontrivial, its order is at most $9\cdot4=36$. If only one is
nontrivial, it is at most $3\cdot8=24$. If neither is nontrivial,
it is at most $8$. Applying the same descent argument to the
joint image proves the last assertion.
\end{proof}

\begin{lemma}[A point bound for supersingular formal groups]
\label{lem:B-RD31-point}
Let $F$ be a one-dimensional formal group over a complete DVR
of mixed characteristic $(0,3)$, of absolute ramification $a$,
whose special fibre has height at least two. If $F$ has a
rational point of exact order $3^n$, then
\[
                         a\geq3^n.                     \tag{B-RD31.2}
\]
\end{lemma}
\begin{proof}
Normalize the valuation to give a uniformizer value $1$.
Write $[3](T)=3T+\sum_{j\geq2}b_jT^j$. For an invariant
differential $u(T)dT$, with $u(T)$ an integral unit series, the
identity $u([3](T))[3]'(T)=3u(T)$ shows that $j b_j$ is divisible by $3$.
Thus $b_2\in3\mathcal O$. Since the special fibre has height
at least two, $b_3,\ldots,b_8$ lie in the maximal ideal.
For $x$ of positive integral valuation $v$ this gives
\[
 v([3]x)\geq\min\{a+v,\,1+3v,\,9v\}.
                                                        \tag{B-RD31.3}
\]
A nonzero point $y$ of order $3$ satisfies
$v(y)\leq(a-1)/2$. Indeed, if $v(y)>(a-1)/2$, the linear term
$3y$ has strictly smaller valuation than all terms of degree
$2$ through $8$. It also has smaller valuation than all terms
of degree at least $9$: if $a\geq2$, then
$8v(y)>4(a-1)\geq a$, with strict inequality also at $a=2$;
if $a=1$, the integral inequality $v(y)\geq1$ gives it
directly. This would contradict $[3]y=0$.

For a point of exact order $3^n$, write
$v_i=v([3]^i x)$ for $0\leq i<n$. The final point satisfies
$v_{n-1}\leq(a-1)/2$. Each earlier $v_i$ also lies below this
bound: (B-RD31.3) strictly increases every positive valuation,
so a value already above that bound could never reach the
final one. Consequently $a\geq2v_i+1$, and
(B-RD31.3) gives $v_{i+1}\geq3v_i+1$.
Starting from $v_0\geq1$, induction yields
\[
 v_{n-1}\geq\frac{3^n-1}{2}.
\]
Combining the two inequalities gives $a\geq3^n$.
\end{proof}

\begin{proposition}[The mixed formal and filtered case is detected by $27$-torsion]
\label{prop:B-RD31-mixed}
Suppose $p=3$, $e_K=2$, and $\mu_3\subset K$.
Let $M$ be a height-two potentially good group with supersingular
special fibre. Let $N$ admit a filtration with height-one graded
pieces and with sum of their dimensions equal to $1$.
Then
\[
                         M[27]\not\simeq N[27].
                                                        \tag{B-RD31.4}
\]
\end{proposition}
\begin{proof}
Pass to the completed maximal unramified extension, which
preserves $e_K=2$ and the given torsion comparison. A potentially
good height-one character is then $\chi_3^b\epsilon$, where
$b\in\{0,1\}$ and $\epsilon$ has values in $\{1,-1\}$.
This is the height-one classification used in
Proposition~\ref{prop:B-RD5}. The determinant of $T_3M$ is
$\chi_3\epsilon_0$, with $\epsilon_0$ likewise quadratic;
this follows by taking the determinant in crystalline comparison,
as in Proposition~\ref{prop:B-RD18}.

An asserted isomorphism transports the first height-one
subgroup of $N[27]$ to a cyclic direct summand $C\subset M[27]$
whose character is $\chi_3^b\epsilon$. If $b=0$, twist $M$ by
$\epsilon$ to obtain a rational point of order $27$. If $b=1$,
form the BT quotient $M'=M/C$. The image of $M[27]$ in $M'[27]$
is cyclic of order $27$ and has character
$\epsilon_0\epsilon^{-1}$. Twisting $M'$ by this quadratic
character again gives a rational point of exact order $27$.
The quotient and twist remain potentially good, one-dimensional,
and supersingular. A finite flat closure of the kernel in a good
model justifies the assertion for the quotient; the twist becomes
trivial over a quadratic extension.

By Lemma~\ref{lem:B-RD31-finite}, the resulting group acquires
good formal reduction over an extension of ramification degree
at most $12$. Its absolute ramification is therefore at most
$24$. Its rational point extends to a point of this formal
group: finite flat torsion is proper over the DVR, and its
connected special fibre has only the identity geometric point.
Lemma~\ref{lem:B-RD31-point} would require absolute ramification
at least $27$, a contradiction.
\end{proof}

\begin{theorem}[B-RD31: three possibilities in the smallest odd-prime case]
\label{thm:B-RD31}
\status{Proved; two independent complete proof audits and primary-source checks passed.}
In B.4 one has
\[
                         2\leq\nu_{\min}(3,2,2)\leq4.
                                                        \tag{B-RD31.5}
\]
More precisely, $M[243]$ determines the prime-to-$3$ character
of every height-two group satisfying the manuscript monodromy
condition. Any counterexample to determination by $M[27]$ or
$M[81]$ must involve two groups which themselves have
potentially good supersingular reduction. Thus deciding the
two levels $27$ and $81$ in that restricted case decides
$\nu_{\min}(3,2,2)$.
\end{theorem}
\begin{proof}
The lower bound is Proposition~\ref{prop:B-RD29-p3-level-two}.
For the upper bound, let $M[243]\simeq N[243]$ and work after
completed maximal unramified base change. For a height-two
monodromy filtration, let $d\in\{0,1,2\}$ be the sum of the
dimensions of its good graded models. Its determinant character
on inertia is $\chi_3^d\epsilon$, with $\epsilon$ finite
quadratic. The image of $\chi_3$ modulo $9$ has order $3$:
$\mu_3\subset K$, whereas $e_K=2$ excludes $\mu_9\subset K$.
Equality of determinants modulo $9$ therefore forces equal $d$
and equal $\epsilon$ for $M$ and $N$.

If $d=0$, all good graded models are etale. Their special
inertia character is the Brauer lift of the semisimplification
of the generic representation modulo $3$, evaluated on
prime-to-$3$ elements. It is therefore determined already by
$M[3]$. If $d=2$, Cartier duality reduces to $d=0$; the
cyclotomic twist has trivial finite descent inertia. This
also treats extensions between the graded pieces, since
semisimplification and the character are additive.

Suppose $d=1$. Each group either has a filtration by height-one
pieces, or is itself potentially good and supersingular.
Indeed a potentially ordinary height-two graded piece has a
canonical connected--etale filtration on its good model;
uniqueness and descent of this filtration produce the claimed
height-one filtration over $K$.
For two groups with height-one filtrations, the graded generic
characters are $\chi_3^{b_i}\epsilon_i$. Since
$\chi_3\equiv1\pmod3$ on inertia, their finite characters
$\epsilon_i$ and hence their sum are determined by the
semisimplified $3$-torsion. A pair of different kinds is excluded
already at level $27$ by Proposition~\ref{prop:B-RD31-mixed}.

It remains to compare two potentially supersingular groups.
Lemma~\ref{lem:B-RD31-finite} gives a common good formal
extension $L$ with $e_L\leq72$. Since
\[
                         72\leq2\cdot3^4,
\]
Bondarko's Theorem~E(1), with $u=4$ and $\ell=1$, extends the
generic $3^5$-torsion comparison to an isomorphism between the
special $3$-torsion models, equivariantly for descent. The
special Dieudonne modules modulo $3$ consequently have the
same prime-to-$3$ Brauer character. A finite-order matrix of
order prime to $3$ on a free Witt module has eigenvalues which
are Teichmuller lifts of its residual eigenvalues. Thus this is
the required equality of characteristic-zero inertia
characters. The same reductions, without the final application
of Theorem~E, give the assertion about possible counterexamples
at levels $27$ and $81$.
\end{proof}


\begin{remark}[Exact auxiliary certificate and remaining gate]
The program \path{code/verify_height_two_small_bounds.py} checks the
quaternion matrix identities, all $49$ joint-image numerical cases,
$16000$ formal-valuation inequalities and the $36$ determinant
comparisons used above. It passed in $0.05041721599991433$ seconds;
its JSON output is retained. These finite checks certify the
explicit algebra and inequalities, not the existence or
nonexistence of a counterexample at either remaining level.
At this intermediate stage neither $27$ nor $81$ was decided for arbitrary pairs
of potentially supersingular groups over fields of absolute
ramification $2$.
\end{remark}

\paragraph{Subsequent exact determination.}
The remaining small-height alternatives isolated here are settled in
Theorem~\ref{thm:B-RD42}: the exact manuscript index is
$\nu_{\min}(3,2,2)=2$. The intermediate bounds and conditional
reductions above are retained as steps in that proof.
